% !TeX root = ../main.tex

\xchapter{绪论}{Introductions}

\xsection{选题意义及应用背景}{Backgrounds and Significance}

随着无线网络基础设施在家庭、办公与公共空间中的持续普及，利用现有 Wi-Fi 设备开展环境感知和行为识别逐渐成为室内智能感知的重要实现方式。相较于依赖摄像头或可穿戴设备的方案，基于 Wi-Fi 的感知方法无需额外贴身部署，且不直接采集可见图像信息，因此在部署便捷性、使用接受度以及隐私敏感场景适用性方面具有一定优势。在这一背景下，围绕商用无线信号开展室内动作识别研究，已成为无线感知领域值得关注的研究方向。随着 Wi-Fi 在家庭、办公与公共场所的广泛部署\cite{ma2019wifi,oughton2021wifi5g}以及物联网终端数量的持续增长\cite{zeinab2017iot}，利用现有通信基础设施进行无线感知，已成为低部署成本、易扩展的实现路径之一。

在 Wi-Fi 感知中，人体移动、遮挡和肢体动作会改变室内多径传播结构，并进一步反映到无线信道响应的时间变化上。信道状态信息相比接收信号强度指标能够提供更丰富的细粒度描述，不仅包含不同子载波上的响应差异，也能够表征多天线链路上的传播变化，因此更适合用于刻画复杂动作造成的动态扰动\cite{ma2019wifi}。围绕 CSI 的室内动作识别研究已经从早期依赖人工设计特征和传统分类器的方法，逐步发展到利用深度模型自动学习时空表征的方法，识别性能与模型表达能力均得到提升。但在类别不断扩展、旧数据难以重复获得以及标注样本受限等实际条件下，现有方法仍难以直接满足长期更新需求。

与实验室条件下预先给定类别集合的识别任务不同，真实场景中的动作类别往往具有动态扩展特征。系统在初始部署阶段通常只覆盖少量核心动作，后续会随着监测目标变化逐步引入新的行为类型。如果每次类别扩展都重新采集全部历史类别与新增类别数据，并在统一数据集上重新训练模型，不仅会带来较高的时间与计算开销，也会显著增加数据组织与标注成本。对于需要持续运行的 Wi-Fi 感知系统而言，这种更新方式缺乏现实可操作性。因此，使模型能够在保留既有识别能力的同时逐步吸收新类别知识，是该类系统走向长期应用必须解决的关键问题。

此外，Wi-Fi 感知数据的采集过程本身也受到较多条件限制。与图像数据相比，CSI 数据采集通常依赖稳定的收发设备布置、特定场景下的重复录制以及后续预处理与人工标注，整个流程更容易受到场地条件、设备状态和人员配合程度的影响。当识别类别逐步增加时，重新组织完整数据采集往往需要较高的人力成本和较长的实验周期，这使得依赖全量数据反复重训的方式在很多场景中并不现实。因此，从固定类别条件下的离线识别转向面向类别扩展的持续学习，是 Wi-Fi 室内动作识别进一步走向长期应用过程中需要解决的重要问题。

基于 CSI 的人类活动识别（Human Activity Recognition, HAR）已采用长短期记忆（LSTM）、循环神经网络（RNN）与卷积神经网络（CNN）等深度模型从原始或预处理后的 CSI 数据中提取特征并完成分类\cite{ma2019wifi}。然而，这类方法通常假定训练前可获取全部类别数据，模型训练完成后若要识别新增动作类别，往往需要从头重训，带来较长训练时间与较高存储开销。现实场景中，用户行为、环境布局与活动类型会随时间变化，数据往往以流式方式到来，模型需能持续适应新类别，而传统批量学习难以满足流式数据下的更新需求。类增量学习关注的是：当新类别按阶段持续到达时，模型如何在主要利用当前阶段数据完成更新的同时，继续保持对历史类别的统一识别能力。与一次性获得全部类别数据的离线训练不同，这类任务更强调模型在持续更新过程中的知识保留与新类吸收能力；其核心挑战在于灾难性遗忘——在固定容量网络中直接优化新类会覆盖旧类表征，导致旧类识别性能急剧下降\cite{zhou2024cil}。这对应于神经网络中的稳定性—可塑性困境：模型若过于稳定则难以充分学习新类别，若过于可塑则易遗忘旧知识。研究面向 Wi-Fi CSI 的类增量学习方法，有助于在保持对旧动作记忆的同时吸收新动作知识，从而提升室内活动识别系统在动态环境中的适应性。

在此基础上，少样本类增量学习（Few-Shot Class-Incremental Learning, FSCIL）进一步引入了新类监督极度受限的条件，即每个新增类别通常只提供 1 至 5 个标注样本；模型因此需要在样本极少与类别持续扩展同时存在的情况下完成更新，其难点明显高于常规类增量学习\cite{tian2024fscilsurvey}。现实应用中，CSI 数据采集与标注成本较高：高精度 Wi-Fi 接入点与专用采集设备需专业人员安装调试，CSI 数据通常具有规模大、存储开销高和标注依赖人工等特点，而标注质量又会直接影响后续模型效果。正因如此，在健康监护、工业巡检等成本敏感场景中，提高有限样本的利用效率、减少对大规模标注数据的依赖，具有较强现实意义。与传统视觉任务相比，Wi-Fi CSI 场景下的少样本增量问题还具有更强的特殊性。一方面，CSI 本质上是对无线信道传播状态的时序观测结果，其特征不仅与人体动作本身有关，还与空间布局、遮挡关系、多径传播结构以及设备布置方式密切相关；同一动作在不同环境、不同用户甚至不同时间段下都可能表现出较大差异。另一方面，一些细粒度动作在信号层面具有较强相似性，例如幅度接近的肢体摆动、节奏相近的连续动作等，容易在类别持续扩展时带来新旧类别之间的混淆。因此，在 Wi-Fi CSI 条件下开展少样本类增量学习，模型不仅需要在极少样本下形成对新类别的有效表征，还需要尽可能保持旧类别判别边界的稳定性，其难度高于一般的固定类别少样本识别任务。

从应用角度看，这种新增类别样本有限但类别仍会继续增加的情况在无线感知任务中并不少见。在康复监护、养老看护和工业巡检等场景中，新增动作通常具有较强的场景依赖性和阶段性，采集者难以像构建离线数据集那样为每一个新类别准备大量标注样本。若模型不能在少样本条件下较快完成新类学习，系统的后续扩展将受到限制；若在学习新类的过程中旧类性能明显下降，又会影响系统长期运行的稳定性。因此，研究面向 Wi-Fi CSI 的少样本类增量学习，不仅有助于提高模型对新类别的适应能力，也有助于提升系统在长期应用中的可扩展性。

例如，老年人护理场景下新增康复动作类别时，往往难以获取大量标注样本；智能家居中用户偏好变更带来的新动作类型，也需模型在少量样本下完成适配。因此，本文关注的不是固定类别集下的离线识别性能提升，而是 Wi-Fi CSI 室内动作识别在类别持续扩展条件下的持续学习能力；进一步地，本文分别讨论新增类样本相对充足和极少两类典型场景。本论文围绕上述两个方向，分别提出面向 Wi-Fi CSI 的类增量学习框架与少样本类增量学习方案，并在 XRF 数据集上进行验证。

\xsection{国内外研究现状分析}{Current Research Situation}

\xsubsection{面向 Wi-Fi 信号室内动作识别任务发展现状}{Research Status of Indoor Action Recognition Using Wi-Fi Signals}

面向室内场景的人类活动识别研究经历了由可穿戴传感、视觉感知到无线感知不断拓展的发展过程，其中基于 Wi-Fi 信号的动作识别因能够复用现有通信基础设施而受到持续关注。该方向的核心思路在于，将人体活动对无线传播路径造成的扰动转化为可建模的信号变化，再通过特征分析或模式学习实现行为判别。随着商用 Wi-Fi 设备可获取更稳定的信道信息，研究重点也逐步从早期的粗粒度状态检测转向更细粒度的动作识别与行为理解，应用场景则进一步延伸到家庭护理、健康监测、安防感知和智能环境交互等领域\cite{ma2019wifi}。

从技术演进看，早期 Wi-Fi 感知研究更多围绕 RSSI 展开，相关方法主要服务于室内定位或粗粒度状态感知。微软在 2000 年推出的 RADAR 系统\cite{bahl2000radar}通过测量接收信号强度实现室内定位，是该领域的先驱性工作；随后 HORUS 系统\cite{youssef2004horus}与 Bose 等的研究\cite{bose2007pathloss}同样采用 RSSI 进行定位建模。由于 RSSI 仅反映接收功率变化，其对复杂多径环境下细微动作差异的刻画能力有限，难以支撑更精细的行为识别任务。随着研究逐步转向 CSI，利用商用设备获取更高分辨率的信道响应成为可能，人体动作对传播路径造成的扰动也因此能够以更细致的形式被记录下来。以 WiHear、E-eyes 和 WiFall 为代表的一批工作表明\cite{wang2014wihear,wang2014eeyes,wang2016wifall}，基于 CSI 的方法已经能够支持语音、日常活动和跌倒检测等多种任务，为后续室内动作识别研究奠定了基础。早期工作多依赖传统的时频域分析与手工特征设计，随着深度学习技术的成熟，研究者逐步转向端到端学习范式，使得基于 Wi-Fi 的室内活动识别在精度与泛化能力上得到提升。

现有基于 Wi-Fi CSI 的室内动作识别方法大体可以从表征方式上分为经验特征驱动方法、传播机理建模方法和深度表征学习方法。第一类方法通常对 CSI 进行滤波、编码或统计变换，并利用构造得到的分布特征完成匹配或分类，代表性思路是将特定动作对应的信号模式视为可检索的“指纹”；这类方法实现相对直接，但对环境变化较为敏感\cite{wang2014eeyes}。第二类方法尝试从人体运动与信道扰动之间的关系出发建立描述模型，例如分析 CSI 变化速度、活动过程与动作类别之间的对应关系，从而增强识别的可解释性\cite{wang2015csiactivity}。第三类方法则借助卷积网络等深度模型，直接从原始或预处理后的 CSI 数据中学习判别特征，在复杂场景中通常表现出更强的表征能力\cite{zeng2014cnn,cho2018divide}。总体来看，不同方法在实现成本、鲁棒性与精度之间各有取舍，但大多数研究仍建立在固定类别集和离线训练条件之上，对类别逐步扩展时的持续更新问题关注不足。

总体来看，现有研究已经从不同角度验证了商用 Wi-Fi 设备用于室内动作识别的可实施性，也说明该方向在方法设计与任务拓展上具有较丰富的发展路径。此外，鲁勇等对基于 Wi-Fi 信号的人体行为感知技术进行了系统梳理，从定位、检测、识别等多个任务维度总结了 Wi-Fi 感知的发展脉络与关键挑战，进一步说明了 Wi-Fi 感知在人类活动识别中的研究价值\cite{luyong2019wifiharreview}。不过，Wi-Fi 室内动作识别在实际部署中仍受到多方面约束，例如环境变化、设备差异、连续数据到达方式以及存储资源限制等因素，都会削弱模型在跨场景条件下的稳定性与泛化表现\cite{ma2019wifi}。XRF55\cite{wang2024xrf55}等数据集在类别数量（55 类）与样本规模上为增量学习研究提供了数据基础。该数据集涵盖人—物交互、人—人交互、健身、身体运动与人—机交互等五大类室内活动，从日常动作（如走路、坐下）到复杂交互（如转呼啦圈、演奏乐器）均有覆盖，既能反映现实场景的多样性，也满足增量学习对类别数与数据量的要求。然而，现有工作多面向固定类别集的离线训练，对类别持续扩展条件下的模型更新讨论不足。当类别集合随业务需求扩展时，往往缺乏系统化的增量更新方案；而面向 Wi-Fi CSI 的少样本类增量学习研究更为稀缺，既需应对 CSI 信号的高维与时变特性，又需在极少量样本下维持旧类判别能力。在此背景下，将类增量学习与少样本类增量学习的理论与方法引入面向 Wi-Fi CSI 的室内动作识别任务，有助于弥合现有工作在类别动态扩展场景中的不足。

\xsubsection{面向 Wi-Fi 感知的持续学习模式研究现状}{Research Status of Continuous Learning for Wi-Fi Sensing}

将持续学习引入面向 Wi-Fi 的室内活动识别，需要同时应对两方面挑战：一方面，人类活动的多样性与环境的动态性要求系统能够实时适应新出现的活动类型，并保持对已学活动的识别能力；另一方面，Wi-Fi 信号对环境因素敏感（如房间布局、物体摆放、人员移动等），且 CSI 数据的采集与标注成本较高，使得在新类别出现时获取大量标注样本往往不现实\cite{ma2019wifi}。

对于持续到达的动态数据，传统模型往往难以同时兼顾新知识吸收与旧知识保持，学习新任务后旧任务性能明显退化的现象较为常见。这种灾难性遗忘会直接限制 Wi-Fi 感知系统在长期运行中的持续识别能力。从更一般的机器学习视角看，类别增量学习要求模型在不断接收新类别数据的同时保持对旧类别的判别能力，灾难性遗忘、类别不平衡与评价协议差异是该方向的核心问题。朱飞等对类别增量学习的任务设定、代表方法与性能评价进行了系统总结，为理解本文所研究的 Wi-Fi CSI 场景下类增量学习问题提供了通用理论参照\cite{zhufei2023cilreview}。针对 Wi-Fi 场景下的类增量学习，已有研究开始尝试将通用持续学习机制迁移到无线感知任务中。例如，WiCAR\cite{licar2024wicar}通过结合记忆回放、知识蒸馏和类别平衡策略，缓解增量训练中的遗忘与样本失衡问题；WiCI-HAR\cite{ding2023wicihar}则采用分阶段训练方式，在增量更新时减弱特征提取器与分类器之间的相互干扰，并借助注意力改进网络结构和多损失约束提升对新类别样本的利用效果。这些工作表明，面向 Wi-Fi 感知的类增量学习已经具备初步可行性，但现有实验通常仍停留在类别规模较小、增量轮次有限的设置下，对复杂动作区分能力和长期扩展稳定性的讨论仍然不够充分。此外，上述方法均假设每个增量任务可访问充足的新类样本，而在健康监护、工业巡检等场景中，新增动作类别往往难以在短期内采集大量标注数据，这使得少样本条件下的类增量学习成为亟待探索的方向。

在新增类别样本受限的情形下，少样本学习为 Wi-Fi 感知提供了另一条值得关注的技术路径。相关研究通常通过度量学习、图结构建模或注意力机制提高模型对少量支持样本的利用效率。例如，有工作利用图卷积与注意力模块共同建模动作样本之间的关联关系，以增强跨场景条件下的识别能力\cite{zhang2021csigdam}；也有研究将原始 CSI 映射到更具判别性的嵌入空间，并结合基于距离的预测策略完成低样本条件下的活动识别\cite{ding2021oneshot}。这类方法说明，Wi-Fi CSI 在样本有限时仍具有可学习的判别结构，但它们大多面向固定类别体系下的快速适配，并未进一步讨论类别持续增加时旧类保持与新类学习之间的矛盾。因此，少样本学习能够为 Wi-Fi 感知中的数据受限问题提供支持，但尚不能直接替代少样本类增量学习研究。在实际部署条件下，反复组织大规模 CSI 采集并据此重新训练模型往往代价过高，因此仅依赖完全重训的思路缺乏现实操作性。相比之下，面向 Wi-Fi 感知的少样本类增量学习仍处于探索阶段，其应用潜力还有待进一步挖掘。相较于计算机视觉领域已在 CIFAR-100、miniImageNet 等基准上形成较为成熟的 FSCIL 方法体系，面向 Wi-Fi CSI 的少样本类增量学习在任务定义、评价协议与算法适配等方面仍存在研究空白。具体而言，当前的研究缺口可归纳为三点：（1）现有 Wi-Fi CIL 方法多假设每个增量任务可访问充足的新类样本；（2）现有 Wi-Fi 少样本学习多不考虑类别持续扩展的场景；（3）现有视觉领域的 FSCIL 方法难以直接迁移到 Wi-Fi CSI 时序模态，需针对 CSI 时序特性开展专门设计。综上，现有研究尚未同时覆盖 Wi-Fi 室内动作识别中类别持续扩展与新增类样本稀缺这两类典型情形。基于上述问题，本文将研究重点放在 Wi-Fi CSI 条件下的类增量学习与少样本类增量学习两类任务上，并结合 CSI 的时序特点设计对应方法。

基于上述分析可以看出，面向 Wi-Fi CSI 室内动作识别的持续学习研究目前仍处于不断探索阶段，在任务设定与方法设计上尚未形成较为统一的技术路线。现有工作一方面更多围绕新增类样本相对充足的类增量学习情形展开，重点在于缓解灾难性遗忘和新旧类别不平衡；另一方面，与少样本相关的研究又主要集中在固定类别条件下的快速识别，对类别集合持续扩展时模型表征稳定性与泛化能力之间的平衡考虑不足。也就是说，Wi-Fi 感知中的类增量学习和少样本类增量学习虽然都服务于模型的持续更新，但二者所面对的约束条件和关键问题并不相同，需要分别研究。

基于此，本文在后续研究中将这两类问题分开讨论，并按照由一般到困难的顺序展开：首先面向新增类样本相对充足的情形，研究如何通过样本重放、分类器扩展和知识巩固等方式缓解灾难性遗忘，构建适用于 Wi-Fi CSI 的类增量学习方法；随后进一步面向新增类样本极少的情形，研究如何在极少监督条件下兼顾新类学习与旧类保持。这样安排的原因在于，两类问题的约束条件并不相同。先讨论新增类样本相对充足的情形，再进一步讨论极少样本条件下的增量学习，更便于说明本文方法是如何逐步展开的。

\xsection{论文的主要研究内容}{Main Research Contents}

围绕上述两类持续更新场景，本文开展了对应的方法设计与实验研究，主要工作如下：

1）针对类增量学习场景，提出回放—扩展—巩固的学习框架，通过类均值选样记忆回放、冻结骨干下的分类器扩展以及双边知识蒸馏巩固缓解灾难性遗忘。在 XRF 数据集上的实验表明，该方法在末次增量阶段（Session 4）准确率达 73.57\%，优于 iCaRL、BiC、UCIR 以及 SimpleCIL、IL2A 等对比方法。

2）针对少样本类增量学习场景，提出基于时序增强与图注意力的学习方案，通过多尺度时序编码、时序注意力融合、增强图注意力网络与双分类器架构，提升少样本条件下的新类适应能力与旧类保持能力。在 XRF 数据集 5-way 1-shot 设定下的实验表明，该方法末次增量会话（Session 5）准确率达 74.24\%，性能下降率（PD）为 15.76 个百分点，平均准确率（AA）达 78.48\%。

\xsection{论文的组织结构}{Structure of Thesis}

全文共分为五章，整体结构如图~\ref{fig:thesis-structure} 所示，各章主要内容安排如下。

\begin{figure}[H]
  \centering
  \includegraphics[width=0.92\textwidth,height=0.40\textheight,keepaspectratio]{fig1-1-overview.png}
  \caption{本文组织结构与研究内容示意}
  \label{fig:thesis-structure}
\end{figure}

第一章 绪论。在第一章中，首先介绍 Wi-Fi 信号用于室内活动识别的研究背景与应用场景；其次，梳理无线感知与持续学习相关研究现状，并说明类增量学习与少样本类增量学习在该任务中的研究动机；最后概述本文的主要研究内容与章节安排。

第二章 相关理论和知识背景。本章梳理与后续方法设计直接相关的理论基础，依次介绍 Wi-Fi CSI 感知原理、类增量学习、少样本学习以及少样本类增量学习的主流方法与核心挑战。

第三章 面向 Wi-Fi 信号动作识别的类增量学习方法研究。本章针对类增量学习场景，提出回放—扩展—巩固框架，并在 XRF 数据集上与多种基线方法进行比较与消融实验。

第四章 面向 Wi-Fi 信号动作识别的少样本类增量学习方法研究。本章针对少样本类增量学习场景，提出基于时序增强与图注意力的学习方案，并在 XRF 数据集上进行实验验证与消融分析。

第五章 结论与展望。对全文工作进行总结，并结合实验结果与方法局限，说明后续可进一步完善的方向。
